\documentclass[UTF8,zihao=-4,fontset=none,twoside]{ctexart}

\usepackage[a4paper,top=37mm,bottom=35mm,left=28mm,right=26mm,footskip=18mm]{geometry}
\usepackage{fontspec}
\usepackage{xeCJK}
\usepackage{fancyhdr}
\usepackage{needspace}
\usepackage{xcolor}
\usepackage[hidelinks]{hyperref}

\setmainfont{Times New Roman}
\setCJKmainfont{STFangsong}
\setCJKmonofont{STFangsong}
\newCJKfontfamily\titlefont{Songti SC}[BoldFont={Songti SC Bold},AutoFakeBold=2]
\newCJKfontfamily\heiti{STHeiti}
\newCJKfontfamily\kaiti{Kaiti SC}
\newCJKfontfamily\fangsong{STFangsong}
\newCJKfontfamily\pagefont{Songti SC}

\hypersetup{
  pdftitle={AI 工具使用声明},
  pdfauthor={晏妍、王彦翔},
  pdfsubject={海河少女可信 RAG 虚拟主播系统 AI 工具使用说明},
  pdfkeywords={AI 工具使用声明, ChatGPT, Codex, 阿里云百炼, Ollama}
}

\XeTeXlinebreaklocale "zh"
\XeTeXlinebreakskip=0pt plus 1pt minus 0.1pt
\setlength{\parindent}{2em}
\setlength{\parskip}{0pt}
\setlength{\emergencystretch}{1em}
\raggedbottom

\pagestyle{fancy}
\fancyhf{}
\fancyfoot[RO]{{\pagefont\zihao{4}\thepage\hspace*{1em}}}
\fancyfoot[LE]{{\pagefont\zihao{4}\hspace*{1em}\thepage}}
\renewcommand{\headrulewidth}{0pt}
\renewcommand{\footrulewidth}{0pt}
\fancypagestyle{plain}{
  \fancyhf{}
  \fancyfoot[RO]{{\pagefont\zihao{4}\thepage\hspace*{1em}}}
  \fancyfoot[LE]{{\pagefont\zihao{4}\hspace*{1em}\thepage}}
  \renewcommand{\headrulewidth}{0pt}
  \renewcommand{\footrulewidth}{0pt}
}

\newcommand{\maintitle}[1]{%
  \begin{center}
    {\titlefont\bfseries\zihao{2}\fontsize{22pt}{35pt}\selectfont #1\par}
  \end{center}
  \vspace{14pt}
}
\newcommand{\levelone}[1]{%
  \Needspace{76pt}\par\vspace{6pt}
  \noindent{\heiti\zihao{3}\fontsize{16pt}{34pt}\selectfont #1\par}
  \nobreak
}
\newcommand{\leveltwo}[1]{%
  \Needspace{76pt}\par\vspace{2pt}
  \noindent{\kaiti\zihao{3}\fontsize{16pt}{34pt}\selectfont #1\par}
  \nobreak
}
\newcommand{\levelthree}[1]{%
  \Needspace{84pt}\par
  \noindent{\fangsong\zihao{3}\fontsize{16pt}{28pt}\selectfont #1\par}
  \nobreak
}
\newcommand{\bodytext}[1]{%
  {\fangsong\zihao{3}\fontsize{16pt}{28pt}\selectfont #1\par}
}
\newcommand{\infoitem}[2]{%
  \noindent\hspace*{2em}{\fangsong\zihao{3}\fontsize{16pt}{28pt}\selectfont #1：#2\par}
}
\newcommand{\model}[1]{\mbox{#1}}

\begin{document}
\thispagestyle{fancy}

\maintitle{海河少女可信 RAG 虚拟主播系统\\AI 工具使用声明}

\levelone{一、作品基本信息}
\infoitem{作品名称}{海河少女 基于可信 RAG 的天津文化 AI 虚拟主播系统}
\infoitem{申报方向}{方向 A AI 智能体创作系统设计}
\infoitem{学校}{天津师范大学}
\infoitem{作者}{晏妍、王彦翔}
\infoitem{指导教师}{王欣雨}
\infoitem{声明日期}{2026 年 9 月}

\levelone{二、声明目的与适用范围}
\bodytext{本声明用于说明“海河少女”作品从需求分析、软件开发、知识库构建、问答生成、语音播报到论文整理过程中实际使用的 AI 工具、模型及其作用边界。本声明将开发过程中的 AI 协作与作品运行时的 AI 服务分开记录，以便评审者判断团队原创内容、模型生成内容与程序自动处理之间的关系。}

\bodytext{本项目使用 AI 作为方案与文稿协作、软件开发辅助、文本向量计算、受证据约束的文本生成、语义复核和语音合成工具。\mbox{AI 不替代}团队确定文化主题、创作角色形象、选择事实来源、审查最终表达和承担参赛责任。}

\levelone{三、开发与文稿协作中使用的 AI 工具}
\leveltwo{（一）OpenAI ChatGPT}
\bodytext{ChatGPT 用于项目早期方案讨论、功能边界梳理、系统架构的备选方案比较、论文提纲组织以及部分中英文表述辅助。它提供的内容属于建议稿和候选表述，不作为天津历史、非物质文化遗产或作品评测数据的事实来源。团队对采用的文字进行重写、删改和事实核对，论文中的文化信息仍以列明的学术、官方与审核资料为依据。}

\leveltwo{（二）OpenAI Codex \model{gpt-5.6-sol}}
\bodytext{Codex 使用 \model{gpt-5.6-sol} 模型参与项目仓库分析、前后端代码协作、测试脚本与构建命令执行、运行问题定位、LaTeX 论文修改、PDF 导出和版式检查。团队先提供具体需求、设计目标和修改意见，再审查文件差异、运行测试并检查\mbox{最终界面与导出文件}。\mbox{Codex 生成或修改}的代码不因“已生成”而被视为可用，只有经过实际构建、单元测试、浏览器交互或 PDF 逐页检查的结果才进入最终作品。}

\bodytext{Codex 与 ChatGPT 仅用于开发和文稿协作，不是参赛演示时的问答服务。作品运行时没有调用 OpenAI API，也没有将 ChatGPT 或 Codex 输出当作检索证据展示给观众。}

\levelone{四、作品运行时使用的 AI 模型}
\leveltwo{（一）百炼 \model{qwen3.7-text-embedding}}
\bodytext{阿里云百炼 \model{qwen3.7-text-embedding} 用于生产配置下的文本向量化。入库程序将人工审核后的知识块转换为 1024 维向量，查询时再将用户问题转换为同维度向量。向量相似度与\mbox{SQLite FTS5 词法检索}结果经排名融合后，选出\mbox{最多 6 个}候选证据块。该模型只负责表示语义相似性，不直接生成事实回答，也不决定哪些资料可以入库。}

\leveltwo{（二）百炼 \model{qwen3.7-flash-2026-07-15}}
\bodytext{固定快照 \model{qwen3.7-flash-2026-07-15} 承担受约束的回答生成和独立语义复核。生成调用只接收当次检索返回的候选片段，\mbox{并被要求}以结构化 JSON 返回段落、引用标识和拒答状态。生成完成后，系统使用独立调用检查回答段落是否被所引片段完整支持。这一复核仍属于工程防线，不是对模型绝对无错的承诺。}

\bodytext{服务端不直接信任模型返回的引用。程序会校验引用标识是否来自当次候选集合，并检查回答中的阿拉伯数字和日期是否出现在对应证据中。任一主题、证据、引用、数字日期或语义检查失败时，系统返回固定拒答，而不要求模型继续补写。}

\leveltwo{（三）百炼 \model{qwen3-tts-flash-2025-11-27} 与 Momo 音色}
\bodytext{百炼 \model{qwen3-tts-flash-2025-11-27} 用于将已经通过全部校验的中文回答转换为语音。系统使用模型提供的 Momo 中文音色，通过非实时 HTTP 调用生成\mbox{24 kHz 单声道 WAV 音频}。本项目没有进行真人声音克隆，也没有收集参赛者或观众的声纹数据。}

\bodytext{语音接口不接收前端任意输入的文本。只有通过回答校验的文字才会由服务端签发一次性语音令牌，令牌过期或使用后立即失效。因此，语音合成是证据校验之后的表达环节，不参与文化事实判断。}

\leveltwo{（四）Ollama 本地开发模型}
\bodytext{开发阶段使用 Ollama 运行两类本地模型。}
\bodytext{\model{qwen3-embedding:4b} 用于验证本地向量生成、索引构建和混合检索流程；\model{qwen3.5:27b} 用于验证离线回答生成、结构化输出、引用约束和拒答逻辑。}
\bodytext{本地模型便于在不消耗云端调用的情况下反复调试，但与百炼生产配置的模型、向量维度和延迟不同。论文中的生产评测结果以百炼配置为准，不将本地基线结果写作云端模型表现。}

\levelone{五、AI 工具在项目流程中的使用情况}
\leveltwo{（一）需求分析与方案设计}
\bodytext{ChatGPT 与 Codex 辅助梳理功能清单、比较 RAG 检索方案、分析风险和整理实施步骤。项目的天津文化主题、“海河少女”角色定位、可信问答的核心目标、展示范围和取舍均由团队确定。}

\leveltwo{（二）视觉形象与角色制作}
\bodytext{“海河少女”三维角色、人物设定、服装纹理、发型配色、角色草图和相关视觉素材由参赛团队原创绘制与制作。}
\Needspace{112pt}
\bodytext{本作品未使用 AI 绘图工具生成角色立绘、三维模型、服装纹理或宣传图，也未使用 AI 视频生成服务制作演示画面。}
\Needspace{84pt}
\bodytext{VRoid Studio、Three.js 和 VRM 运行库属于人物制作或图形渲染软件，不列为 AI 内容生成工具。}

\leveltwo{（三）资料整理与知识库构建}
\bodytext{团队决定资料是否入库，并核对标题、作者或发布单位、页码、网址和使用边界。程序完成文档解析、切块、哈希记录和索引构建，向量模型只负责计算语义表示。对扫描文献的文本实行按页人工校对，角色世界观与历史事实分类标记，避免把艺术创作设定表述为历史结论。}

\leveltwo{（四）问答、引用与语音播报}
\bodytext{系统先进行主题、实时性和提示词攻击检查，再执行词法与向量两路检索。\mbox{FTS5 负责}词法匹配，向量模型负责语义召回。文本生成模型只能使用当次返回的候选片段，完成生成后还必须通过引用白名单、数字日期一致性和独立语义复核。只有通过全部检查的文字才能进入 TTS；证据不足、主题越界或检查失败时，系统明确拒答且不播报。}

\leveltwo{（五）测试、论文与交付材料}
\bodytext{Codex 辅助执行单元测试、检索评测、前端构建、浏览器检查和 PDF 渲染检查。ChatGPT 与 Codex 参与论文结构、措辞和版式协作，但测试数据、模型名称、知识库规模、引用来源和最终结论由团队对照程序输出与原始材料确认。论文不将 AI 生成的建议性文字列为学术引用，也不用 AI 自报的评价代替可复现测试。}

\levelone{六、人工审核与责任边界}
\levelthree{1. 参赛团队负责确定作品主题、角色定位、视觉风格、文化符号和功能边界。}
\levelthree{2. 参赛团队负责判断资料是否可信与可用，核对原文、页码、网址和引用片段。}
\levelthree{3. 参赛团队负责审查 AI 协作产生的代码和文字，并通过测试、差异检查和视觉验收决定是否采用。}
\levelthree{4. 参赛团队负责确认程序评测数据与最终文稿一致，不将局部测试结果扩大为对任意场景的保证。}
\levelthree{5. 参赛团队对作品原创性、知识性内容、代码可用性、数据安全和最终提交文件承担责任。}

\bodytext{在实际问答中，模型输出还受到程序规则约束。问答主题超出限定知识库、引用标识不属于当次检索结果、数字或日期缺少证据、语义复核不通过时，系统均不输出事实性回答。这些程序检查可以降低无依据生成的概率，但不替代团队对资料和演示内容的最终审查。}

\levelone{七、数据、隐私与敏感信息处理}
\bodytext{作品演示采用文字输入，未接入麦克风和语音识别，不建立用户账户，不收集观众身份信息、人脸、声纹或定位数据，不将对话内容设计为持久化用户历史。运行时不进行网页搜索，只将当次问题与回答所必需的少量审核片段发送给已配置的百炼模型服务。}

\bodytext{百炼 API Key、业务空间标识、主机名和专用服务地址只保存在本机已忽略的 \texttt{.env} 配置中，不写入源码、知识库、论文、截图、视频或参赛压缩包。对外记录仅保留模型名称、向量维数、调用结果与不含凭据的验收信息。}

\levelone{八、原创性与工具边界说明}
\bodytext{角色形象、角色设定、三维模型、服装纹理、发型配色、视觉草图及相关展示素材均由参赛团队原创绘制和制作。AI 不是该角色视觉成果的作者，也未替代团队进行文化符号选择、造型调整和最终审美判断。文档编辑软件、代码库、图形渲染引擎、排版系统和浏览器自动化工具不因参与工程流程而被等同为\mbox{AI 内容生成工具}。}

\bodytext{知识性内容的事实责任仍由原始资料和参赛团队承担。模型只能在已给定证据中组织回答，引用卡片中展示的标题、发布者、页码、片段和网址用于回到原始依据。\mbox{AI 输出本身}不作为其自身正确性的证明。}

\levelone{九、评测记录与能力边界}
\bodytext{2026 年 9 月 13 日，项目使用百炼 1024 维生产向量完成固定 60 题评测：40 个库内问题全部在 Top-6 候选中命中预标来源，20 个越界、歧义与提示词攻击问题全部被拒答。正式 API 连续问答 20/20 符合预期，其中 2 轮真实语音请求返回 HTTP 200，后端 30 项测试全部通过。}
\Needspace{112pt}
\bodytext{这些数据用于说明当前语料、题集和配置下的工程表现，不代表大语言模型在理论上绝对不会产生幻觉。}

\bodytext{本作品的能力边界由限定知识库决定。对未入库的主题、需要实时查询的信息、证据存在冲突或语义复核无法通过的问题，系统应当拒答。如后续更换模型、更新资料或调整切块规则，必须重建索引并重新运行固定评测，不直接沿用旧结论。}

\levelone{十、最终声明}
\bodytext{参赛团队确认，本文档已列明开发协作与作品运行中实际使用的 AI 工具及其主要用途。团队未将 AI 生成内容伪装为未经协助的人工成果，也未将模型输出当作无需审核的文化事实。作品的创意决策、原创视觉资产、资料准入、代码验收、事实复核和最终提交均由参赛团队负责。}

\end{document}
